\problem{P23}{越界前后的两个间隙}{更新密度}
$U_i$ 独立服从 $\operatorname{Unif}(0,1)$，记 $S_k=\sum_{i=1}^kU_i$，$S_0=0$，$N=\min\{k:S_k>1\}$。令 $A=1-S_{N-1}$ 为越界前距离，$O=S_N-1$ 为越界量。求 $(A,O)$ 的联合密度、最后一步长度 $J=U_N$ 的密度，以及 $\Prob(O>A)$。
\solution 对 $0<s<1$，第 $k$ 个正累计和的密度为 $s^{k-1}/(k-1)!$。所有正累计和的密度相加得到
\[
u(s)=\sum_{k\ge1}\frac{s^{k-1}}{(k-1)!}=e^s.
\]
第一步不可能超过 $1$，故最后越界前的累计和必为正，没有来自 $S_0$ 的原子。若越界前在 $s$，下一步为 $v$，则联合密度贡献为 $e^s\dd s\dd v$，条件是 $s+v>1$。由于增量为正，这个跨越步骤唯一。

变换 $a=1-s$、$o=s+v-1$ 的雅可比绝对值为 $1$，所以
\[
\boxed{f_{A,O}(a,o)=e^{1-a},\quad 0<a<1,\quad0<o<1-a.}
\]
积分 $\int_0^1e^{1-a}(1-a)\dd a=1$，确为概率密度。最后一步 $J=A+O$，因此
\[
\boxed{f_J(j)=\int_0^je^{1-a}\dd a=e-e^{1-j},\quad0<j<1.}
\]
要有 $O>A$，必须 $a<1/2$，故
\[
\boxed{\Prob(O>A)=\int_0^{1/2}e^{1-a}(1-2a)\dd a
=2\sqrt e-e.}
\]
\insight 把所有可能的越界前累计和相加，再接上最后一个增量，就能得到越界前后两个量的联合分布，而无须分别枚举停止次数。
